\section{Structure--Transport Evidence and Minimum Comparability}

Structure--transport reasoning often compresses a multistep chain into ``denser means more conductive.'' The audited abstracts do support study-specific covariation, but they also reveal why the shortcut is unsafe. A hot-pressing study reports density and the fraction of resistance associated with grain boundaries as pressing temperature changes \cite{david2015microstructure}. A submicron Nb-containing study reports phase, relative density, microstructure, total conductivity, and activation energy together \cite{yan2020submicron}. A crucible comparison reports composition, density, conductivity, and response to air exposure \cite{xia2016ionic}. Each case involves more than a scalar density.

Relative density also does not uniquely encode pore location, grain size, boundary chemistry, surface state, or spatial uniformity. The abstract on uniform densification explicitly couples microscopic and macroscopic observations with heating and support conditions \cite{guo2023uniform}. The solution/solid-state comparison attributes transport differences within its specimens to microstructural features while also varying route and Al inclusion \cite{kosir2022comparative}. We retain those as authors' abstract-reported interpretations, not as a universal density--conductivity law.

Transport labels must remain explicit. Total conductivity includes the response resolved under the study's equivalent-circuit treatment; bulk and grain-boundary terms are not interchangeable with it. Test temperature, electrode configuration, specimen geometry, frequency range, fitting choice, and activation-energy convention can change the reported observable. The cover-powder abstract is useful precisely because it identifies total conductivity and activation energy alongside density and grain-boundary composition \cite{zhang2019densification}; it does not make those values directly comparable to a record that omits temperature or reports a different partition.

We therefore apply the minimum record in a fixed sequence:

\begin{enumerate}
\item Establish composition and provenance: nominal formula, measured chemistry when available, dopant and additive, and whether support comes from title, abstract, excerpt, or checked full text.
\item Establish material history: precursors, mixing and powder route, lithium excess, calcination, shaping, and every sintering field relevant to pressure, time, ramp, atmosphere, cover, and crucible.
\item Establish specimen state: phase method, density basis, grain and boundary descriptors, porosity or spatial sampling, and exposure history.
\item Establish transport meaning: total/bulk/grain-boundary quantity, temperature, electrodes, analysis convention, replicates, and uncertainty.
\end{enumerate}

The output is categorical, not a score. If relevant fields align, two observations may be directly compared. If a known field differs, they remain parallel qualitative evidence that can motivate a controlled test. If a decisive field is NR, the comparison is insufficiently specified. This avoids two symmetrical errors: treating all nominally similar LLZO as identical, and treating every difference as evidence for the processing variable named in a title.

Contribution B is thus more than a reporting checklist. It is the decision layer attached to Contribution A's process map. Contribution C then uses the same fields to specify why an unresolved comparison remains unresolved. For example, a low-temperature additive case cannot discriminate density mediation from boundary-phase mediation unless density and boundary response are separately measured. A route comparison cannot isolate precursor homogeneity unless composition and thermal budget are matched. These are test-design consequences of the evidence map, not claims that the relevant experiments are absent from the field.
